% Part: first-order-logic
% Chapter: tableaux
% Section: derivations

\documentclass[../../../include/open-logic-section]{subfiles}

\begin{document}

\iftag{FOL}
      {\olfileid{fol}{tab}{der}}
      {\olfileid{pl}{tab}{der}}

\olsection{\usetoken{P}{tableau}}

\begin{explain}
قد بيّنا الافتراض وقواعد الاستدلال؛ ومن هذين الأصلين تُبنى
!!^{tableau}s استقرائيًا. فكل !!{tableau} إما فرع واحد قوامه
افتراض أو افتراضات، وإما حاصل من !!a{tableau} سابق بإجراء
إحدى قواعد الاستدلال على أحد فروعه.
\end{explain}

\begin{defn}[!!^{tableau}]
ال!!^a{tableau} للافتراضات~$\sFmla{S_1}{!A_1}$، \dots،
$\sFmla{S_n}{!A_n}$، حيث كل~$S_i$ إما~$\True$ وإما~$\False$،
شجرة منتهية من ال!!{signed formula}s تستوفي الشرطين الآتيين:
\begin{enumerate}
\item تكون ال!!{signed formula}s العليا وعددها~$n$ في الشجرة هي
  $\sFmla{S_i}{!A_i}$، واحدة تحت الأخرى.
\item تنتج كل !!{signed formula} في الشجرة ليست من الافتراضات عن تطبيق
  صحيح لقاعدة استدلال على !!a{signed formula} في الفرع فوقها.
\end{enumerate}
وفرع !!a{tableau} \emph{مغلق} إذا وفقط إذا اجتمع فيه
$\sFmla{\True}{!A}$ و$\sFmla{\False}{!A}$، وإلا فهو \emph{مفتوح}.
وال!!^a{tableau} الذي أغلقت فروعه جميعًا هو \emph{!!{tableau} مغلق}
لمجموعة افتراضاته. وأما !!a{tableau} الذي لم يغلق، أي بقي له فرع
مفتوح واحد على الأقل، فيسمّى \emph{مفتوحًا}.
\end{defn}

\begin{ex}
  مجموعة الافتراضات وحدها !!a{tableau}، ولا يلزم أن تكون مغلقة؛
  بل لا تغلق إلا إذا اشتملت منذ البدء على زوج من ال!!{signed formula}s
  $\sFmla{\True}{!A}$ و$\sFmla{\False}{!A}$.

  ومن !!a{tableau}، مفتوحًا كان أو مغلقًا، نحصل على آخر أكبر منه
  إذا أجرينا قاعدة استدلال على !!a{signed formula}
  $\sFmla{S}{!A}$ فيه. فتُلحق القاعدة !!{signed formula}s، واحدة
  أو أكثر، بنهاية كل فرع يمر بالورود المقصود من~$\sFmla{S}{!A}$.


  خذ مثلًا الافتراض~$\sFmla{\True}{!A \land \lnot !A}$.
  فال!!{tableau} المؤلّف منه وحده مفتوح، وصورته:
  \begin{center}
    \begin{tableau}{}
      [\sFmla{\True}{\formula{A} \land \lnot \formula{A}}, just=\TAss]
    \end{tableau}{}
  \end{center}
  نجعل منه !!{tableau} جديدًا بإجراء~$\TRule{\True}{\land}$
  على الافتراض؛ فتُلحق بال!!{tableau} سطران،
  هما~$\sFmla{\True}{!A}$ و$\sFmla{\True}{\lnot !A}$:
  \begin{center}
    \begin{tableau}{}
      [\sFmla{\True}{\formula{A} \land \lnot \formula{A}}, just=\TAss
        [\sFmla{\True}{\formula{A}}, just={\TRule{\True}{\land}[1]},
          [\sFmla{\True}{\lnot \formula{A}}, just={\TRule{\True}{\land}[1]}
          ]
        ]
      ]
    \end{tableau}{}
  \end{center}
  وإذا كتبنا ال!!{tableau}s أثبتنا القواعد المستعملة في الجانب الأيمن.
  فمثلًا~$\TRule{\True}{\land} 1$ يدل على أن ال!!{signed formula}
  في ذلك السطر أُخذت بقاعدة~$\TRule{\True}{\land}$ من
  ال!!{signed formula} في السطر~$1$.
  وقد زيدت في هذا ال!!{tableau} !!{signed formula}s، ولكن واحدة
  منها فقط، وهي~$\sFmla{\True}{\lnot !A}$، تقبل الآن قاعدة،
  أعني~$\TRule{\True}{\lnot}$. وبإجرائها يغلق ال!!{tableau} هكذا:
  \begin{center}
    \begin{tableau}{}
      [\sFmla{\True}{\formula{A} \land \lnot \formula{A}}, just=\TAss
        [\sFmla{\True}{\formula{A}}, just={\TRule{\True}{\land}[1]}
          [\sFmla{\True}{\lnot \formula{A}}, just={\TRule{\True}{\land}[1]}
            [\sFmla{\False}{\formula{A}}, just={\TRule{\True}{\lnot}[3]}, close]
          ]
        ]
      ]
    \end{tableau}{}
  \end{center}
\end{ex}

\end{document}
